\documentclass[a4paper,12pt,twoside]{memoir}

% Castellano
\usepackage[spanish,es-tabla]{babel}
\selectlanguage{spanish}
\usepackage[utf8]{inputenc}
\usepackage[T1]{fontenc}
\usepackage{lmodern} % Scalable font
\usepackage{microtype}
\usepackage{placeins}
\usepackage{makecell}

\RequirePackage{booktabs}
\RequirePackage[table]{xcolor}
\RequirePackage{xtab}
\RequirePackage{multirow}

% Links
\PassOptionsToPackage{hyphens}{url}\usepackage[colorlinks]{hyperref}
\hypersetup{
	allcolors = {red}
}

% Ecuaciones
\usepackage{amsmath}

% Rutas de fichero / paquete
\newcommand{\ruta}[1]{{\sffamily #1}}

% Párrafos
\nonzeroparskip

% Huérfanas y viudas
\widowpenalty100000
\clubpenalty100000

% Imágenes

% Comando para insertar una imagen en un lugar concreto.
% Los parámetros son:
% 1 --> Ruta absoluta/relativa de la figura
% 2 --> Texto a pie de figura
% 3 --> Tamaño en tanto por uno relativo al ancho de página
\usepackage{graphicx}
\newcommand{\imagen}[3]{
	\begin{figure}[!h]
		\centering
		\includegraphics[width=#3\textwidth]{#1}
		\caption{#2}\label{fig:#1}
	\end{figure}
	\FloatBarrier
}

% Comando para insertar una imagen sin posición.
% Los parámetros son:
% 1 --> Ruta absoluta/relativa de la figura
% 2 --> Texto a pie de figura
% 3 --> Tamaño en tanto por uno relativo al ancho de página
\newcommand{\imagenflotante}[3]{
	\begin{figure}
		\centering
		\includegraphics[width=#3\textwidth]{#1}
		\caption{#2}\label{fig:#1}
	\end{figure}
}

% El comando \figura nos permite insertar figuras comodamente, y utilizando
% siempre el mismo formato. Los parametros son:
% 1 --> Porcentaje del ancho de página que ocupará la figura (de 0 a 1)
% 2 --> Fichero de la imagen
% 3 --> Texto a pie de imagen
% 4 --> Etiqueta (label) para referencias
% 5 --> Opciones que queramos pasarle al \includegraphics
% 6 --> Opciones de posicionamiento a pasarle a \begin{figure}
\newcommand{\figuraConPosicion}[6]{%
  \setlength{\anchoFloat}{#1\textwidth}%
  \addtolength{\anchoFloat}{-4\fboxsep}%
  \setlength{\anchoFigura}{\anchoFloat}%
  \begin{figure}[#6]
    \begin{center}%
      \Ovalbox{%
        \begin{minipage}{\anchoFloat}%
          \begin{center}%
            \includegraphics[width=\anchoFigura,#5]{#2}%
            \caption{#3}%
            \label{#4}%
          \end{center}%
        \end{minipage}
      }%
    \end{center}%
  \end{figure}%
}

%
% Comando para incluir imágenes en formato apaisado (sin marco).
\newcommand{\figuraApaisadaSinMarco}[5]{%
  \begin{figure}%
    \begin{center}%
    \includegraphics[angle=90,height=#1\textheight,#5]{#2}%
    \caption{#3}%
    \label{#4}%
    \end{center}%
  \end{figure}%
}
% Para las tablas
\newcommand{\otoprule}{\midrule [\heavyrulewidth]}
%
% Nuevo comando para tablas pequeñas (menos de una página).
\newcommand{\tablaSmall}[5]{%
 \begin{table}
  \begin{center}
   \rowcolors {2}{gray!35}{}
   \begin{tabular}{#2}
    \toprule
    #4
    \otoprule
    #5
    \bottomrule
   \end{tabular}
   \caption{#1}
   \label{tabla:#3}
  \end{center}
 \end{table}
}

%
% Nuevo comando para tablas pequeñas (menos de una página).
\newcommand{\tablaSmallSinColores}[5]{%
 \begin{table}[H]
  \begin{center}
   \begin{tabular}{#2}
    \toprule
    #4
    \otoprule
    #5
    \bottomrule
   \end{tabular}
   \caption{#1}
   \label{tabla:#3}
  \end{center}
 \end{table}
}

\newcommand{\tablaApaisadaSmall}[5]{%
\begin{landscape}
  \begin{table}
   \begin{center}
    \rowcolors {2}{gray!35}{}
    \begin{tabular}{#2}
     \toprule
     #4
     \otoprule
     #5
     \bottomrule
    \end{tabular}
    \caption{#1}
    \label{tabla:#3}
   \end{center}
  \end{table}
\end{landscape}
}

%
% Nuevo comando para tablas grandes con cabecera y filas alternas coloreadas en gris.
\newcommand{\tabla}[6]{%
  \begin{center}
    \tablefirsthead{
      \toprule
      #5
      \otoprule
    }
    \tablehead{
      \multicolumn{#3}{l}{\small\sl continúa desde la página anterior}\\
      \toprule
      #5
      \otoprule
    }
    \tabletail{
      \hline
      \multicolumn{#3}{r}{\small\sl continúa en la página siguiente}\\
    }
    \tablelasttail{
      \hline
    }
    \bottomcaption{#1}
    \rowcolors {2}{gray!35}{}
    \begin{xtabular}{#2}
      #6
      \bottomrule
    \end{xtabular}
    \label{tabla:#4}
  \end{center}
}

%
% Nuevo comando para tablas grandes con cabecera.
\newcommand{\tablaSinColores}[6]{%
  \begin{center}
    \tablefirsthead{
      \toprule
      #5
      \otoprule
    }
    \tablehead{
      \multicolumn{#3}{l}{\small\sl continúa desde la página anterior}\\
      \toprule
      #5
      \otoprule
    }
    \tabletail{
      \hline
      \multicolumn{#3}{r}{\small\sl continúa en la página siguiente}\\
    }
    \tablelasttail{
      \hline
    }
    \bottomcaption{#1}
    \begin{xtabular}{#2}
      #6
      \bottomrule
    \end{xtabular}
    \label{tabla:#4}
  \end{center}
}

%
% Nuevo comando para tablas grandes sin cabecera.
\newcommand{\tablaSinCabecera}[5]{%
  \begin{center}
    \tablefirsthead{
      \toprule
    }
    \tablehead{
      \multicolumn{#3}{l}{\small\sl continúa desde la página anterior}\\
      \hline
    }
    \tabletail{
      \hline
      \multicolumn{#3}{r}{\small\sl continúa en la página siguiente}\\
    }
    \tablelasttail{
      \hline
    }
    \bottomcaption{#1}
  \begin{xtabular}{#2}
    #5
   \bottomrule
  \end{xtabular}
  \label{tabla:#4}
  \end{center}
}



\definecolor{cgoLight}{HTML}{EEEEEE}
\definecolor{cgoExtralight}{HTML}{FFFFFF}

%
% Nuevo comando para tablas grandes sin cabecera.
\newcommand{\tablaSinCabeceraConBandas}[5]{%
  \begin{center}
    \tablefirsthead{
      \toprule
    }
    \tablehead{
      \multicolumn{#3}{l}{\small\sl continúa desde la página anterior}\\
      \hline
    }
    \tabletail{
      \hline
      \multicolumn{#3}{r}{\small\sl continúa en la página siguiente}\\
    }
    \tablelasttail{
      \hline
    }
    \bottomcaption{#1}
    \rowcolors[]{1}{cgoExtralight}{cgoLight}

  \begin{xtabular}{#2}
    #5
   \bottomrule
  \end{xtabular}
  \label{tabla:#4}
  \end{center}
}



\graphicspath{ {./img/} }

% Capítulos
\chapterstyle{bianchi}
\newcommand{\capitulo}[2]{
	\setcounter{chapter}{#1}
	\setcounter{section}{0}
	\setcounter{figure}{0}
	\setcounter{table}{0}
	\chapter*{#2}
	\addcontentsline{toc}{chapter}{#2}
	\markboth{#2}{#2}
}

% Apéndices
\renewcommand{\appendixname}{Apéndice}
\renewcommand*\cftappendixname{\appendixname}

\newcommand{\apendice}[1]{
	%\renewcommand{\thechapter}{A}
	\chapter{#1}
}

\renewcommand*\cftappendixname{\appendixname\ }

% Formato de portada
\makeatletter
\usepackage{xcolor}
\newcommand{\tutor}[1]{\def\@tutor{#1}}
\newcommand{\course}[1]{\def\@course{#1}}
\definecolor{cpardoBox}{HTML}{E6E6FF}
\def\maketitle{
  \null
  \thispagestyle{empty}
  % Cabecera ----------------
\noindent\includegraphics[width=\textwidth]{cabecera}\vspace{1cm}%
  \vfill
  % Título proyecto y escudo informática ----------------
  \colorbox{cpardoBox}{%
    \begin{minipage}{.8\textwidth}
      \vspace{.5cm}\Large
      \begin{center}
      \textbf{TFG del Grado en Ingeniería Informática}\vspace{.6cm}\\
      \textbf{\LARGE\@title{}}
      \end{center}
      \vspace{.2cm}
    \end{minipage}

  }%
  \hfill\begin{minipage}{.20\textwidth}
    \includegraphics[width=\textwidth]{escudoInfor}
  \end{minipage}
  \vfill
  % Datos de alumno, curso y tutores ------------------
  \begin{center}%
  {%
    \noindent\LARGE
    Presentado por \@author{}\\ 
    en Universidad de Burgos --- \@date{}\\
    Tutor: \@tutor{}\\
  }%
  \end{center}%
  \null
  \cleardoublepage
  }
\makeatother

\newcommand{\nombre}{Licinio Fernández Maurelo} %%% cambio de comando

% Datos de portada	
\title{Evaluación de simulaciones de proyectos en GitHub orientadas a docencia.}
\author{\nombre}
\tutor{Carlos López Nozal}
\date{\today}

\begin{document}

\maketitle


\newpage\null\thispagestyle{empty}\newpage


%%%%%%%%%%%%%%%%%%%%%%%%%%%%%%%%%%%%%%%%%%%%%%%%%%%%%%%%%%%%%%%%%%%%%%%%%%%%%%%%%%%%%%%%
\thispagestyle{empty}


\noindent\includegraphics[width=\textwidth]{cabecera}\vspace{1cm}

\noindent D. nombre tutor, profesor del departamento de nombre departamento, área de nombre área.

\noindent Expone:

\noindent Que el alumno D. \nombre, con DNI dni, ha realizado el Trabajo final de Grado en Ingeniería Informática titulado título de TFG. 

\noindent Y que dicho trabajo ha sido realizado por el alumno bajo la dirección del que suscribe, en virtud de lo cual se autoriza su presentación y defensa.

\begin{center} %\large
En Burgos, {\large \today}
\end{center}

\vfill\vfill\vfill

% Author and supervisor
\begin{minipage}{0.45\textwidth}
\begin{flushleft} %\large
Vº. Bº. del Tutor:\\[2cm]
D. Carlos López Nozal
\end{flushleft}
\end{minipage}
\hfill
%\begin{minipage}{0.45\textwidth}
%\begin{flushleft} %\large
%Vº. Bº. del co-tutor:\\[2cm]
%D. nombre co-tutor
%\end{flushleft}
%\end{minipage}
%\hfill

\vfill

% para casos con solo un tutor comentar lo anterior
% y descomentar lo siguiente
%Vº. Bº. del Tutor:\\[2cm]
%D. nombre tutor


\newpage\null\thispagestyle{empty}\newpage




\frontmatter

% Abstract en castellano
\renewcommand*\abstractname{Resumen}
\begin{abstract}
En este primer apartado se hace una \textbf{breve} presentación del tema que se aborda en el proyecto.
\end{abstract}

\renewcommand*\abstractname{Descriptores}
\begin{abstract}
Palabras separadas por comas que identifiquen el contenido del proyecto Ej: servidor web, buscador de vuelos, android \ldots
\end{abstract}

\clearpage

% Abstract en inglés
\renewcommand*\abstractname{Abstract}
\begin{abstract}
A \textbf{brief} presentation of the topic addressed in the project.
\end{abstract}

\renewcommand*\abstractname{Keywords}
\begin{abstract}
keywords separated by commas.
\end{abstract}

\clearpage

% Indices
\tableofcontents

\clearpage

\listoffigures

\clearpage

\listoftables
\clearpage

\mainmatter
\capitulo{1}{Introducción}

En la actualidad, la enseñanza a distancia goza de gran importancia debido a ventajas inherentes como la deslocalización de profesores y alumnos, la libertad de agenda de los alumnos a la hora de seguir la formación y el uso eficiente de los recursos educativos que permite a los centros de formación ampliar su capacidad de alumnado al no estar sujetos a limitaciones propias de la enseñanza presencial como aulas y laboratorios.

Sin embargo, la enseñanza a distancia de materias relacionadas con la ciencia, la tecnología y  la ingeniería  presenta dificultades derivadas de la propia naturaleza de estas disciplinas: a menudo requieren prácticas en el laboratorio para dotar al alumno de las competencias necesarías . \cite{bourne2019} menciona  cómo, históricamente, algunas de las necesidades especiales de los estudios universitarios de ingeniería no han sido bien cubiertas por los métodos de enseñanza online.

La emergencia de las tecnologías de la información ha permitido el desarrollo de laboratorios virtuales. Su acceso online y el software que los soporta permite a los alumnos una reproducción cada vez más fiel del trabajo realizado en el laboratorio.

Como \cite{tejado2018}  señala, los laboratorios virtuales son esenciales para que los alumnos adquieran conocimientos prácticos que se convertirán en  activos fundamentales en su  carrera profesional. En el desarrollo de laboratorios virtuales, uno de los focos de acción es facilitar la autoevaluación para mejorar la experiencia de aprendizaje de los alumnos y reducir la carga de trabajo de los docentes.Una autoevaluación clara y concisa lleva al alumno a una reflexión sobre sus propios errores, convirtiendo una evaluación formativa en una evaluación formadora \cite{sánchez2009}.

La enseñanza de un modelo agile de gestión de proyectos mediante Github y el sistema de control de versiones Git proporciona el marco de este trabajo de fin de grado. Corresponde a la asignatura Gestión de proyectos - bloque gestión ágil, del tercer curso del grado en ingeniería informática. En esta asignatura se plantean dos actividades a modo de laboratorio virtual:

\begin{itemize}
	\item Gestión de tareas de un proyecto software. El alumno realiza una simulación de planificación de tareas de un proyecto software.
	\item Simulación de control de versiones de un proyecto open source. El alumno simula la gestión de versiones de código utilizando un proyecto disponible en un repositorio público.
\end{itemize}

Los supuestos de ambas actividadades proporcionan el  detalle paso por paso de la simulación a realizar además de la referencia al repositorio original. 

\imagen{planning-simu-showcase.png}{Gestión de tareas de un proyecto software: extracto de información para la simulación}{1}

El producto entegable de cada una de las prácticas debe constar de diversas capturas de pantalla que constaten la fidelidad de la simulación con respecto al original.Para la evaluación de cada una de las prácticas, existe una rúbrica de evaluación que actualmente es aplicada de forma manual en base al producto entregado.

\imagen{planning-simu-rubric.png}{Gestión de tareas de un proyecto software: extracto de rúbrica para la evaluación}{1}

La evaluación de ambas actividades puede automatizarse en gran medida. La información de las tareas de un proyecto y del control de versiones de Github es accesible vía API. A partir del contraste de la información del proyecto referencia y el proyecto de la simulación podemos dar contestación a gran parte de los elementos de la rúbrica. Del mismo modo, podemos explotar esta comparación para proporcionar retroalimentación al alumno de aquellos elementos de la simulación que concuerdan o discrepan con el original.
\capitulo{2}{Objetivos del proyecto}

La realización de este proyecto tiene como objetivo principal el desarrollo de un software que permita automatizar la evaluación y retroalimentación de los siguientes casos de estudio de simulación planteados como laboratorio virtual en la asignatura gestión de proyectos del grado en ingeniería informática:

\begin{enumerate}
	\item Simulación de gestión de tareas de un proyecto software. 
	\item Simulación de control de versiones de un proyecto open source.
\end{enumerate}

Para acometer lo anterior el software desarrollado deberá implementar las siguientes funcionalidades:

\begin{enumerate}
	\item Obtención de información en un intervalo temporal de la gestión de tareas de un proyecto alojado en GitHub con la extensión ZenHub activada:
	\begin{enumerate}
		\item Sprints.
		\begin{enumerate}
			\item Título. 
			\item Fecha de finalización.
			\item Estado.	
		\end{enumerate}	
		\item Tareas.	
		\begin{enumerate}
			\item Descripción. 
			\item Etiquetas.
			\item Puntos de historia.
			\item Comentarios.
			\item Estado.	
		\end{enumerate}	
	\end{enumerate}		
	\item Obtención de información en un intervalo temporal del control de versiones de un proyecto alojado en GitHub:
	\begin{enumerate}
		\item Commits. 
			\begin{enumerate}
				\item Identificador único (SHA, Simple Hashing Algorithm). 
				\item Fecha.
				\item Autor.
			\end{enumerate}		
	\end{enumerate}
	\item Comparación de la información de la gestión de tareas de dos proyectos.
	\item Comparación de la información del control de versiones de dos proyectos.			
	\item Aplicación de rúbrica de evalución  sobre:
	\begin{enumerate}
		\item La información de gestión de tareas y del control de versiones de un repositorio. 
		\item La comparación de la información de gestión de tareas y de control de versiones de dos repositorios, el original y el resultante de la simulación.	
	\end{enumerate}			  
\end{enumerate}

Los objetivos  técnicos del desarrollo son:

\begin{enumerate}
	\item Uso de técnicas de gestión de proyectos ágiles.
	\item Cumplimiento de estándares de calidad de software y aseguramiento de calidad mediante análisis automatizados con Sonarqube.
	\item Integración con GitHub. 
	\item Integración con ZenHub.
	\item Uso de una arquitectura hexagonal \cite{hexagonal-architecture} que desacople la funcionalidad desarrollada de los orígenes de datos. Esto facilitará la extensibilidad de la funcionalidad simplificando la  integración con otros gestores de tareas y gestores de repositorios Git.
\end{enumerate}


\capitulo{3}{Conceptos teóricos}

\section{Conceptos de la gestión de proyectos ágiles utilizados en el laboratorio virtual.}

\subsection{Historia de usuario.}

Es una descripción informal de una funcionalidad del sistema a desarrollar. Se documentan desde la perspectiva del usuario final del sistema.

TBD - Incluir captura historia de usuario.

\subsection{Sprint/Milestone/Hito.}

Es un periodo de tiempo durante el cual un trabajo específico, definido en las distintas tareas que componen el sprint. En el laboratorio virtual es de relevancia la siguiente información:

TBD - Incluir captura tarea

\begin{enumerate}
	\item Fecha de inicio.
	\item Fecha de fin.
	\item Puntos de historia de usuario: total de puntos de historia de usuario de las tareas incluidas en el sprint.	
\end{enumerate}

\subsection{Tarea/Issue.}

Cada historia de usuario es dividida en unidades de trabajo individuales. En el laboratorio virtual se plantea especificar la siguiente información:

\begin{enumerate}
	\item Descripción de la tarea.
	\item Miembro del equipo al que se le asigna la tarea.
	\item Puntos de historia de usuario.
	\item Tipo de tarea. TBD - Indicar tipos
	\item Etiquetas. TBD - Indicar cuales	
\end{enumerate}

TBD - Incluir captura tarea

En aquellos proyectos que necesiten para su comprensión y desarrollo de unos conceptos teóricos de una determinada materia o de un determinado dominio de conocimiento, debe existir un apartado que sintetice dichos conceptos.

Algunos conceptos teóricos de \LaTeX \footnote{Créditos a los proyectos de Álvaro López Cantero: Configurador de Presupuestos y Roberto Izquierdo Amo: PLQuiz}.

\section{Secciones}

Las secciones se incluyen con el comando section.

\subsection{Subsecciones}

Además de secciones tenemos subsecciones.

\subsubsection{Subsubsecciones}

Y subsecciones. 


\section{Referencias}

Las referencias se incluyen en el texto usando cite \cite{wiki:latex}. Para citar webs, artículos o libros \cite{koza92}.


\section{Imágenes}

Se pueden incluir imágenes con los comandos standard de \LaTeX, pero esta plantilla dispone de comandos propios como por ejemplo el siguiente:

\imagen{add_remote.png}{Autómata para una expresión vacía}{1}

\section{Listas de items}

Existen tres posibilidades:

\begin{itemize}
	\item primer item.
	\item segundo item.
\end{itemize}

\begin{enumerate}
	\item primer item.
	\item segundo item.
\end{enumerate}

\begin{description}
	\item[Primer item] más información sobre el primer item.
	\item[Segundo item] más información sobre el segundo item.
\end{description}
	
\begin{itemize}
\item 
\end{itemize}

\section{Tablas}

Igualmente se pueden usar los comandos específicos de \LaTeX o bien usar alguno de los comandos de la plantilla.

\tablaSmall{Herramientas y tecnologías utilizadas en cada parte del proyecto corta}{l c}{herramientasportipodeusocorta}
{ \multicolumn{1}{l}{Herramientas} & App AngularJS  \\}{ 
 \makecell{HTML5 More H \\More}& X \\
} 

\tablaSmall{Herramientas y tecnologías utilizadas en cada parte del proyecto}{l c c c c}{herramientasportipodeuso}
{ \multicolumn{1}{l}{Herramientas} & App AngularJS & API REST & BD & Memoria \\}{ 
HTML5 & X & & &\\
CSS3 & X & & &\\
BOOTSTRAP & X & & &\\
JavaScript & X & & &\\
AngularJS & X & & &\\
Bower & X & & &\\
PHP & & X & &\\
Karma + Jasmine & X & & &\\
Slim framework & & X & &\\
Idiorm & & X & &\\
Composer & & X & &\\
JSON & X & X & &\\
PhpStorm & X & X & &\\
MySQL & & & X &\\
PhpMyAdmin & & & X &\\
Git + BitBucket & X & X & X & X\\
Mik\TeX{} & & & & X\\
\TeX{}Maker & & & & X\\
Astah & & & & X\\
Balsamiq Mockups & X & & &\\
VersionOne & X & X & X & X\\
} 

\capitulo{4}{Técnicas y herramientas}

Esta parte de la memoria tiene como objetivo presentar las técnicas metodológicas y las herramientas de desarrollo que se han utilizado para llevar a cabo el proyecto. Si se han estudiado diferentes alternativas de metodologías, herramientas, bibliotecas se puede hacer un resumen de los aspectos más destacados de cada alternativa, incluyendo comparativas entre las distintas opciones y una justificación de las elecciones realizadas. 
No se pretende que este apartado se convierta en un capítulo de un libro dedicado a cada una de las alternativas, sino comentar los aspectos más destacados de cada opción, con un repaso somero a los fundamentos esenciales y referencias bibliográficas para que el lector pueda ampliar su conocimiento sobre el tema.



\capitulo{5}{Aspectos relevantes del desarrollo del proyecto}

Este apartado pretende recoger los aspectos más interesantes del desarrollo del proyecto, comentados por los autores del mismo.
Debe incluir desde la exposición del ciclo de vida utilizado, hasta los detalles de mayor relevancia de las fases de análisis, diseño e implementación.
Se busca que no sea una mera operación de copiar y pegar diagramas y extractos del código fuente, sino que realmente se justifiquen los caminos de solución que se han tomado, especialmente aquellos que no sean triviales.
Puede ser el lugar más adecuado para documentar los aspectos más interesantes del diseño y de la implementación, con un mayor hincapié en aspectos tales como el tipo de arquitectura elegido, los índices de las tablas de la base de datos, normalización y desnormalización, distribución en ficheros3, reglas de negocio dentro de las bases de datos (EDVHV GH GDWRV DFWLYDV), aspectos de desarrollo relacionados con el WWW...
Este apartado, debe convertirse en el resumen de la experiencia práctica del proyecto, y por sí mismo justifica que la memoria se convierta en un documento útil, fuente de referencia para los autores, los tutores y futuros alumnos.

\capitulo{6}{Trabajos relacionados}

Este apartado sería parecido a un estado del arte de una tesis o tesina. En un trabajo final grado no parece obligada su presencia, aunque se puede dejar a juicio del tutor el incluir un pequeño resumen comentado de los trabajos y proyectos ya realizados en el campo del proyecto en curso. 

\capitulo{7}{Conclusiones y Líneas de trabajo futuras}


TBD Conclusiones.

Se plantean las siguientes líneas de trabajo futuras:

\begin{enumerate}
	\item Ampliar la información a evaluar en las simulaciones . Esto conllevará ampliar el modelo de dominio así como los datos obtenidos en las integraciones con GitHub y ZenHub:
	\begin{enumerate}
		\item Wiki.
		\item Pull requests.
		\item Estimación de tipo póker.
	\end{enumerate} 	
\end{enumerate} 



\bibliographystyle{plain}
\bibliography{bibliografia}

\end{document}
